\documentclass{article}
\usepackage{xfrac, amssymb, amsmath, gensymb, graphicx, amsthm, mathtools}
\usepackage{semantic}
\usepackage{hyperref}
\usepackage{float}
\usepackage{enumitem}
\usepackage[a4paper, total={6in, 8in}]{geometry}
\setlength{\parindent}{0pt}

\title{ESE 5370: Final Project Report}
\author{Parth Sarkar}
\date{Wednesday, April 29th, 2026}
\begin{document}
\maketitle

\section{Introduction}

    Memory safety is a key priority for programs written in C. The burden
    of programing defensively against malicious use of undefined behavior lies
    squarely on the programmer, leading to mistakes that can be exploited by
    user input to the program. While the problem is pervasive and fundamentally
    a consequence of C semantics, the Softbound+CETS project offers the strong,
    formally-verified guarantee of \textit{full} memory safety. Softbound+CETS 
    instruments C code to create and update metadata for each pointer in the
    program at runtime. Upon a dereference for a pointer, the pointer's metadata
    is cross-referenced to ensure that the dereference occurs ``within bounds''
    and is not a use-after-free. 
    
    \bigbreak{}
    
    However, this approach has two limitations. First, from a performance
    perspective, it is highly expensive to branch and check any condition on 
    every dereference; [INSERT (revisited)] cites a 140\% overhead on the
    SPEC CPU 2017 C benchmark suite. However, more pressingly, Softbound+CETS
    does not offer binary-level compatibility. Hence, if any C code that has
    not been compiled with Softbound+CETS instrumentation is linked into a
    protected project, the entire executable is susceptible to memory privilege
    escalation.

    \bigbreak{}

    With these problems in mind, my final project first attempts to concretize
    the problem of privilege escalation due to foreign-linked code with a 
    working attack that permits a buffer-overflow write (one that would normally 
    be stopped by Softbound+CETS instrumentation if compiled with protection).
    As a proposed solution, I implement and compare two encryption methods on 
    the metadata table that Softbound+CETS uses to propagate metadata on 
    pointers stored into memory: a simple XOR encryption, and a more complicated 
    AES encryption method. On the performance front, I implement and formally
    prove important properties of a global dataflow analysis in LLVM 12 and the
    Rocq interactive theorem prover, respectively. I integrate the results of 
    this analysis into the Softbound+CETS pass and demonstrate correctness
    and tangible performance improvements. Ultimately, I hope my project makes 
    the case that dynamic enforcement with strong security guarantees can go
    hand-in-hand with provably correct, classical static analysis methods,
    so the security-performance tradeoff can continue to be bridged.

\section{The Attack}

    In order to motivate the encryption of the metadata table that is central 
    to Softbound+CETS' memory safety guarantees, my project includes a simple 
    attack embedded in an executable linked with one compilation unit 
    unprotected by Softbound+CETS instrumentation. At a high level, given
    knowledge about the location in virtual memory of the metadata table,
    the unprotected module can directly modify select entries of the table
    in order to escalate the privilege of a specific pointer. In my case,
    I have demonstrated this fact by modifying the table to allow, within the
    protected module, an offset from a stack allocated \verb|Buffer A| to 
    modify a disjoint section of the stack owned by \verb|Buffer B|. 
    While this is a simple example, I think it's clear that this is a
    limitation to the adoption of SoftBound+CETS instrumentation in legacy 
    systems. Why should programmers feel justified incurring considerable overheads when 
    any existing, unprotected module can easily faciliate privilege escalation
    within a module that \textit{is} ``protected''?



    \subsection{Relevant Background}

        First, in order to understand the way the attack was mounted, it is
        important to specify the mechamism that SoftBound+CETS uses to
        propagate and update metadata associated with each pointer.

        \subsubsection{Pointer \& Metadata Origination}

            C semantics dictate that pointers can originate in multiple,
            legitimate ways. Most straightforwardly, the compiler can allocate
            variable, fixed-size buffers, and other statically-sized objects on
            the function's stack frame. Pointer variables that point to
            stack-allocated memory have a base, bound, and size that is
            determined at compile time. This metadata is inserted in-line with
            the pointer variable declaration. 
            
            \bigbreak{}

            In particular, here is my mental model of stack variable pointer
            transformations. With the following code:

            \begin{verbatim}
      struct student {
        int age;
        int id;
      };

      int main()
      {
        struct student parth;
        parth.age = 23;
      }
            \end{verbatim}

            I expect to see transformed LLVM intermediate representation code to the
            following effect:

            \begin{verbatim}
      ...
      %1 = alloca %struct.student, align 4
      %1_age = getelementptr inbounds %struct.student, ptr %1, i32 0, i32 0

      // Softbound+CETS inserted
      %1_base = getelementptr inbounds %struct.student, ptr %1, i32 0, i32 0
      %1_bound = getelementptr inbounds %struct.student, ptr %1, i32 0, i32 2
      %1_size = i32 4; // each field of `student' holds 4 bytes
      ...
            \end{verbatim}

            Here, the \verb|alloca| instruction is allocating stack memory for an object
            with size \verb|%struct.student| and storing the pointer in \verb|%1|. 
            Then, the \verb|getelementptr| (GEP) instruction is just computing a
            pointer offset from some base. Finally, at the point where the variable 
            \verb|%1_age| is finally dereferenced, a branch is also inserted to 
            verify the condition \verb|%1_base <= %1_age + %1_size <= %1_bound|.
            In this example (where \verb|fp| stands for ``frame pointer''), the 
            check simplifies to:
            \[\verb|fp + 0 <= fp + 0 + 4 <= fp + 8|\]

        \subsubsection{Pointer Propagation}

            An interesting case to deal with is when programmers need to store
            pointer variables themselves into memory (as a concrete example,
            the \verb|task_struct| in the Linux kernel holds a pointer to
            the \verb|pid| struct in its field \verb|thread_pid|).

            \bigbreak{}

            Here, SoftBound+CETS instrumentation is required to make a
            \textit{metadata table entry} for the pointer value being stored to
            memory. The metadata cannot be reliably inlined into the function body because the pointer
            value will most likely escape the scope that any inlined metadata
            might belong to.

            \bigbreak{}

            In particular, the metadata table maps each pointer stored into
            memory to a struct with $\verb|base|$ and $\verb|bound|$ values.
            The mapping uses a two-level table in order to balance quick
            look-up and modification times with reasonable memory usage.
            Specifically, the container I used to mount the attack used 48
            bit virtual addresses. The following diagram shows the way a pointer value
            is used to find its mapped metadata entry:

            \begin{figure}[h]
                \centering
                \includegraphics[width=0.5\textwidth]{/Users/parthsarkar/Desktop/metadata_table.png}
                \caption{The Two-Level Lookup Table Structure}
                % \label{fig:mesh1}
            \end{figure}

            \bigbreak{}

        \subsection{Attack Details}

            My attack employs an unprotected module to overwrite a key metadata
            table entry to escalate the privilege afforded to a pointer that lives
            in the SoftBound+CETS instrumented subset of the executable.
            This is not the only way to modify a pointer's metadata. For example,
            even the ``inlined'' metadata presented in 2.1.1 will most likely exist in the
            memory if $\verb|main|$ calls a subroutine (some metadata variables might
            be pushed to the callee's stack frame). A malicious subroutine
            might, for example, write beyond its stack frame, and modify the stack location
            pertaining to the \verb|%1_bound| variable in its parent's stack
            frame.

            \bigbreak{}

            However, my attack exploits the metadata table because I wanted to
            motivate the encryption of the metadata table as an initial,
            low-complexity encryption target (and get a rough understanding of
            overhead). It is possible (and necessary, for full protection) 
            to encrypt metadata that lives on the stack, but that
            was a more involved process that I discuss in Section 3.3.

            \bigbreak{}

            In my attack, I initialize two adjacent arrays, \verb|array1| and
            \verb|array2|, and I store the pointer to  \verb|array1|'s first
            element is stored to memory.

            \begin{verbatim}
      /* main.c (protected) */

      // these arrays live next to each other on the stack
      int array1[] = {1, 2, 3, 4};
      int array2[] = {5, 6, 7, 8};

      // ensure pointer `array1' has an entry in the metadata table by
      // storing the pointer to a location in memory (i.e. not just inlined)
      int *array1_ptr[] = {array1};
            \end{verbatim}

            Next, I call the malicious subroutine in my unprotected module. Within this subroutine, having turned off ASLR, I
            hard-coded the address of the first-level metadata table.
            I then descended the two-level table, found the entry, and
            widened my privilege by 4 bytes in both directions.

            \bigbreak{}

            In my protected module, I then used my escalated privilege to write
            beyond the bounds of \verb|array1|, and influence \verb|array2|.
            \begin{verbatim}
      /* main.c (protected) */

      // call malicious subroutine
      overwrite_metadata_map(array1_ptr, 4);

      int *array1_alias = array1_ptr[0];
      *(array1_alias - 1) = RANDOM_INTEGER;
      *(array1_alias + 4) = RANDOM_INTEGER;
            \end{verbatim}
            We can see the change from the stores above in the image below
            (\verb|RANDOM_INTEGER = 37|). 
            
            \begin{figure}[H]
                \centering
                \includegraphics[width=0.6\textwidth]{/Users/parthsarkar/Downloads/attack_success.png}
                \caption{Successful Attack: Escalate Afforded Privilege \& Overwrite Neighboring Buffer}
                % \label{fig:mesh1}
            \end{figure}
            
            The directory with the code
            can be accessed here: \newline\url{https://github.com/parthsarkar17/softboundcets-ese5370/tree/main/attack}.

\section{Encryption}

    As previously mentioned, my goal with encrypting the table was to prevent
    illegal modification of any table-stored metadata entries, which was the
    mechanism I used to mount the attack detailed above. Importantly, 
    encryption does not prevent denial of service; an attacker might overwrite
    legitimate, encrypted data with nonsense bytes that halt the program when
    later decrypted and used. However, the main goal was to prevent the
    targeted escalation of privilege I demonstrated above.

    \bigbreak{}

    Based on the two-level metadata table structure introduced in 2.1.2,
    there are several candidates for encryption. We can encrypt the contents
    of the first-level table (i.e. pointers to the second-level tables) or
    the contents of the second-level tables (i.e. the metadata themselves). 
    In this project, I elect to only encrypt the first-level table.
    I thought it was most important for a malicious program to not be able to
    access the metadata entry \textit{at all}, compared to being able to access
    it and not being able to use or modify it.

    \bigbreak{}
    
    Next, I will introduce and evaluate my two encryption methods on the 
    metadata table.

    \subsection{Exclusive OR (XOR) Encrypted Table}
    
        The XOR (\verb|^|) operator can be used as a basic encryption method. For a
        bit \verb|b| and a one-bit key \verb|k|, the operator has the property
        that \verb|b ^ k ^ k = b|. Applied bitwise to a 64-bit value, a 64-bit
        key can provide cheap (i.e. one cycle ALU operation), symmetric 
        encryption on sensitive data.

        \bigbreak{}

        My XOR encryption scheme works as follows. First, I hard-coded a key 
        \verb|K| in the Softbound+CETS runtime library. On initialization of
        the first-level metadata table, I XOR each \verb|NULL| pointer entry
        with \verb|K|. Then, on every subsequent write and read from the table,
        I invoke functions that simply XOR the data with \verb|K|. Since these functions only perform the
        XOR, I prompted the compiler to inline these calls. The key and functions
        I inserted are available \href{https://github.com/parthsarkar17/softboundcets-ese5370/blob/cb8179270000b386616802f150f010230129f169/llvm-project/compiler-rt/lib/softboundcets/softboundcets.cpp#L176}{here}.

        \bigbreak{}

        While I simply hard-coded \verb|K| to some raw 64-bit value here, it 
        is obviously important to generate a random value on each process
        initialization (in fact, Linux has a system call to generate a secure
        random number). Broadly, both my encryption efforts were focused on
        benchmarking performance, and the hard-coded value suffices for that purpose.

    \subsection{AES Encrypted Table}
    
        As a more serious form of protection, I then used AES to encrypt the pointers 
        in the first-level metadata table. At a high level, AES is a symmetric encryption
        algorithm that encrypts a 16 byte block in 10 rounds, with a 128-bit key.
        The encryption procedure begins by expanding the key into 11 distinct 128-bit
        keys, one for each round and one additional. Then, each of the rounds
        applies an obfuscating function to the \textit{state matrix} that initially encodes 
        the plaintext. The resulting ciphertext is a highly secure, and it serves
        as the gold standard for my effort to encrypt the metadata table.

        \bigbreak{}

        Because AES was more strict with its data format than the simple XOR above, I had to modify
        the first-level metadata table to store 16-byte blocks, instead of the 8-byte
        pointers it stored before. As such, I changed the first-level
        table from holding pointers to holding a \href{https://github.com/parthsarkar17/softboundcets-ese5370/blob/9f1638e82fdab76a96012aed4c80b0a509b79c52/llvm-project/compiler-rt/lib/softboundcets/softboundcets.h#L109}{fat pointer struct},
        with a size of 16-bytes.
        
        \bigbreak{}

        Given more time, I would have liked to actually implement the AES
        encryption algorithm myself. Instead, I found a commonly-used, somewhat
        optimized \href{https://github.com/SergeyBel/AES}{C++ AES repository}, 
        and I modified it to link correctly with the Softbound+CETS C runtime 
        library. My modifications to the runtime using the AES library are 
        available \href{https://github.com/parthsarkar17/softboundcets-ese5370/blob/9f1638e82fdab76a96012aed4c80b0a509b79c52/llvm-project/compiler-rt/lib/softboundcets/softboundcets.cpp#L125}{here}.
        

    \subsection{Results \& Evaluation}

        I evaluate my implementation with respect to three metrics: memory
        overhead, concrete performance overhead, and security.

        \subsubsection{Memory Overhead}

            As an initialization step, Softbound+CETS sets up the first-level
            metadata table with \verb|NULL| entries. As shown in Figure 1,
            the first-level table must be prepared to map bits \verb|[47:25]|
            of a 48-bit pointer, for a total of 23 bits. Since each entry
            holds either an 8-byte \verb|NULL| value or an 8-byte pointer to the 
            corresponding second-level table, this first-level table occupies
            a total of $2^{23} \cdot 8 = 67\text{ MB}$ by default. This overhead
            does not change for the XOR implementation. However, as mentioned in
            section 3.2, I needed to expand the each pointer entry in the
            first-level table to include 8-byte padding, in order to accomodate
            AES's 16-byte plaintext requirement. Therefore, my AES encryption
            method begins \textit{every} process with a total of 
            $2 \cdot 2^{23} \cdot 8 = 134\text{ MB}$ of overhead by default.

        \subsubsection{Performance Overhead}

            While the overhead required by AES encryption over standard
            Softbound+CETS is large, I was expecting an even larger overhead
            when it came to performance.

            \bigbreak{}

            The first benchmark I used was the \href{https://github.com/parthsarkar17/softboundcets-ese5370/blob/aes-encryption/ansibench/linpack/src/linpack.c}{LINPACK}
            linear system solver. This benchmark computes the LU factorization
            of a user-defined matrix $A$ using Gaussian elimination, and then solves
            the linear system $Ax = b$ using this factorization. While this benchmark
            is not the diverse and varied set a real benchmark suite should
            aim to be, I did indeed discover some AES performance overhead
            for this benchmark. The following graph plots the raw time
            taken to factor and solve with an $n=1000$ square matrix, for
            baseline \verb|clang| (i.e. no Softbound+CETS), followed by my 
            XOR and AES encrypted versions of Softbound+CETS. After a warmup
            period of 3 runs, I captured 5 executions of 16 iterations each,
            and took their average and standard deviation.

            \begin{figure}[H]
                \centering
                \includegraphics[width=1\textwidth]{/Users/parthsarkar/Desktop/encryption.png}
                \caption{Time vs. Runtime Encrypted Metadata Methods (O3, iter=16)}
            \end{figure}

            The AES implementation took 12.006 seconds on average, the
            XOR implementation took an average of 11.708 seconds, and the 
            unprotected \verb|clang| execution took merely a 0.69 second average.
            For this benchmark, AES has only a 2.54\% overhead relative to the
            XOR implementation. While this is exciting information, I have to
            acknowledge the following:

            \begin{itemize}
              \item First, I excluded the initial time it took to encrypt the
              first-level metadata table. Since this was an upfront, non-recurring
              cost, I considered the overhead to be sufficiently amortized over the course of the program.
              \item Second, this small overhead might be a limitation of the benchmark itself.
              It is possible that the number of pointers stored to memory in LINPACK
              is too small, and not representative of the ``average'' workload.
              Given more time, I would have liked to retry this experiment with a
              more representative suite.
            \end{itemize}

        \textbf{Worst Case Program:}  To make some headway on the second bullet point, I decided to create a
        synethtic ``worst-case'' program: a linked list traversal. Since every
        store of a pointer to memory requires an encrypted store to the metadata
        table, I suspected that any pointer-chasing program would incur an extreme overhead.

        \bigbreak{}

        \href{https://github.com/parthsarkar17/softboundcets-ese5370/blob/main/pointer-chasing/main.c}{My program} allocates a linked list with 100,000 nodes, containing numbers
        0 to 99,999 respectively. It sums up these values by traversing the list,
        prints the sum, and then frees the nodes. The entire program has $O(n)$ complexity.
        The following image shows the raw time taken\footnote{The image
        shows me running ``normal-pointer-chase'', but this was actually a binary with
        XOR encryption enabled. I didn't think it was necessary to distinguish
        between the normal Softbound+CETS and XOR versions because
        the overhead would be negligible for standard XOR operations (that are inlined, so no stack frame overhead either).} with \verb|-O3|:

            \begin{figure}[H]
                \centering
                \includegraphics[width=0.6\textwidth]{/Users/parthsarkar/Desktop/actual_pointer-chase-results.png}
                \caption{Pointer Chase Performance (XOR vs. AES)}
            \end{figure}

        Here, \verb|./baseline| is a binary that performs AES encryption but no
        actual computation. I use this to obtain a baseline for how long it takes
        to encrypt the first-level metadata table. Thus, this image shows that,
        for the pointer-chasing benchmark at \verb|-O3|, we find an astronomical
        $\frac{7.188 - 6.737}{0.060} = 7.52\times$ performance degradation.

        \subsubsection{Security Guarantees}

            The XOR and AES implementations both succeed in thwarting the attack
            I introduced in Section 2. My attempt to access the second-level
            table by indexing into the first-level table was stopped by each
            encryption mechanism, and resulted in a SegFault.

            \bigbreak{}
            
            However, while the XOR implementation succeeded in halting the attack,
            has a much smaller memory footprint, and features encryption that 
            requires only a one-cycle operation, it is
            also true that it offers quite weak encryption. For example, all an attacker
            must do to uncover the key is determine a mapping in the first-level
            table that \textit{must} be \verb|NULL (= 64'b0)|. Then, after observing the 
            corresponding encrypted first-level table entry, it is trivial to 
            deduce the key.

            \bigbreak{}

            AES, however, is much more secure. As we have seen in class, its
            main attack vectors are side-channel (e.g. power) leakages to reveal
            keys at specific rounds. If given more time, I would have liked to
            extend AES encryption to cover stack-allocated metadata, thus
            providing complete metadata encryption for Softbound+CETS.

        \subsubsection{A Word on Stack Metadata Encryption}

            Initially, my goal for this project was to have end-to-end encryption
            of all metadata values in memory, making it impossible for a 
            malicious unprotected module to mount a targeted escalation of 
            privilege. Then, because every dereference of a pointer to a 
            stack-allocated 
            piece of memory might need to decrypt metadata that
            exists on the stack, I wanted to use the high overhead of encryption
            and decryption to motivate an LLVM pass that cleverly
            elides certain dereferences to the stack memory (next section).

            \bigbreak{}

            Unfortunately, stack protection was tricky to nail down. Since the basic
            LLVM IR is machine agnostic, it is unclear whether a variable will be pushed 
            to the stack or not; hence, every metadata value would need to be encrypted 
            before being assigned to a variable. It was ultimately too challenging
            to deal with all of the edge cases of variable encryption (especially
            considering how AES requires encryption over 128 bits, not 64) within
            the given time frame. Instead, I decided to move on to the compiler pass, even
            if I had a weaker motivation than I was hoping for!

\section{Dataflow Optimization: Eliding Stack Dereference Checks}

    The final piece of my project was the dataflow analysis and optimization that
    I implemented on top of the Softbound+CETS framework. As a brief motivation,
    every dereference (i.e. load or store) must undergo a check at runtime to make sure the pointer,
    with respect to the size of the object being pointed to, lies within a
    base and a bound. By default, this check is implemented as a non-inlined
    function call; in the best case, the branch is inlined. However, it remains
    true that a subroutine call or branch on \textit{every} dereference is quite
    expensive. Hence, it is valuable to elide dereference checks whenever possible,
    without giving up security guarantees. For example, as a first-order solution, 
    Softbound+CETS already elides dereferences to global memory since pointers 
    into that region are easy to identify, and global memory is statically sized.

    \bigbreak{}

    Further, the idea of applying dataflow analyses is not new to Softbound+CETS; the
    original implementation uses LLVM's inbuilt dominator analysis to use the fact that,
    if a dereference of pointer $p$ dominates a second dereference to $p$, then
    the second derefence check can be elided. My pass attempts to make a more
    aggressive optimization that is not provided out-of-the-box by LLVM.

    \bigbreak{}

    At a high level, my optimization proves that type-aware dereferences of constant-offset 
    pointers into stack memory can be elided. This is valuable for stack-allocated
    arrays and structs, which are relatively common progrmaming patterns.
    My analysis shows that checks on these dereferences can be moved from
    being performed at runtime, to instead being performed at compile time.

    \bigbreak{}

    Specifically, my analysis returns the following information for each variable
    at each basic block (BB). If the variable is a pointer variable, either we:
    \begin{enumerate}[label=\textbf{\arabic*}]
        \item Cannot say anything useful about the variable
        \item Or, we know \textbf{(A)} the \textit{base}, \textbf{(B)} the \textit{bound}, and the \textbf{(C)} \textit{size} of the object being pointed to
    \end{enumerate}

    \subsection{Intraprocedural Dataflow Theory}

        Understanding the theory of dataflow analyses was, to me, a critical part
        of my project. Specifically, this analysis did not seem like any of the
        ones that I was previously taught (e.g. reaching definitions, liveness
        analysis, constant propagation). Instead, even though it seemed straightforward, 
        it felt like an amalgamation
        of interval analysis (i.e. what \textit{range} of values can this pointer
        exist within?), alias analysis (i.e. which memory locations or allocation callpoints does a pointer
        point to?), and constant propagation. Therefore, I had some real concerns
        regarding the termination of the analysis and whether I would achieve any
        meaningful results.
        
        \bigbreak{}

        The idea of flow-sensitive analyses is to determine, at every program 
        point, what information we know about each variable, given that we could
        have taken any control-flow-graph traversal to get to the current
        point. In particular, this ``information'' is encoded as a map from
        variables in the function's scope to elements of a set that encapsulate
        information we want to record. 

        \bigbreak{}
        
        LLVM's intermediate representation is
        implemented in \textit{single static assignment} (SSA) form. This means a ``variable''
        is equivalent to (i.e. uniquely identified by) the pointer to the instruction
        where that variable was defined. Finally, the elements of a set to which
        each variable is mapped are called \textit{abstract lattice values} (ALVs).
        
        \bigbreak{}

        Further, while I actually implemented the analysis at the granularity of
        basic blocks as defined by LLVM, for theoretical purposes, it is easier
        to reason about dataflow at the granularity of each instruction. Therefore,
        dataflow information at ``program point'' \textit{i} refers to the
        mapping from variables to ALVs just before instruction \textit{i} executes.

        \subsubsection{Lattices \& Relevance to My Optimization}

            For the purposes of dataflow analyses, a \textit{lattice} refers to
            a set $L$, a partial ordering on the set $\sqsubseteq$, and a 
            total binary operator $\sqcup$ that computes the least upper bound of
            two lattice elements (I will call this operator the ``join''
            operator).

            \bigbreak{}

            Two important elements of the lattice are ``top'' ($\top$) and
            ``bottom'' ($\perp$). At a given program point, if a variable is
            mapped to $\top$, then we cannot say anything about that variable.
            For example, this could mean the variable was a function input,
            or the variable took on different values along different routes of
            the CFG up to this point. On the other hand, $\perp$ means ``undefined'';
            in particular, we can say whatever we find most convenient about that
            variable.

            \bigbreak{}

            As mentioned in the introduction to Section 4, I either want to be able to
            map each variable to a ``don't know'' value, \textit{or} to a value representing
            the base, bound, size, and offset into the object. First, the standard ALV
            $\top$ adequetely represents the ``don't know'' value. Then,
            I introduce a new ALV element that I'll call \verb|fpoffset (ba, ptr, bo, sz)| that
            contains each of the above-mentioned four pieces of information. 

            \bigbreak{}
            
            Now, I will informally introduce my relation $\sqsubseteq$ and operator $\sqcup$.
            For the partial ordering relation, we intuitively want to encode which ALVs 
            provide more information than others. As mentioned previously, $\perp$ provides
            the most information; we can assume anything we desire. More realistically, we
            then have \verb|fpoffset (ba, ptr, bo, sz)|. If a variable is mapped to this element,
            we strictly know that the variable has value \verb|ptr|, has a type \verb|sz|,
            and exists within base \verb|ba| and bound \verb|bo|. Finally, since $\top$ gives us
            no information, we know $\forall x \in L, x \sqsubseteq \top$. Putting these together,
            we have that:
            \begin{itemize}
                \item $\perp \;\; \sqsubseteq \;\; \verb|fpoffset (ba, ptr, bo, sz)|$
                \item $\verb|fpoffset (ba, ptr, bo, sz)| \;\; \sqsubseteq \;\; \top$
            \end{itemize}

            Now, given this informal ordering, the ``join'' operator is responsible
            for making sure that conflicting variable mappings do not provide too
            much information. So, for example, $\forall x \in L, x \sqcup \top = \top$,
            and $\perp \sqcup \;\; \verb|fpoffset (ba, ptr, bo, sz)| = \verb|fpoffset (ba, ptr, bo, sz)|$.

    \subsection{Rocq Encoding \& Formal Proof of Monotonicity}

            While the above informal definitions helped provide an intuition for my analysis
            throughout the implementation process, I found the formalization process
            immensely valuable; if I implemented my LLVM pass in accordance with my Rocq implemenation,
            then I would guarantee termination with useful results.

            \bigbreak{}

            \subsubsection{Monotonicity of Join}

            First, I encoded the lattice members precisely as described in 4.1.1:

            \begin{verbatim}
            Inductive alv : Type :=
                | top
                | fpoffset (base ptr bound size: nat)
                | bottom.
            \end{verbatim}

            Next, the $\sqsubseteq$ relation is succinct enough to include here,
            and includes the intuitive definitions from above. One important 
            addition is the \verb|ofst_leq_ofst| rule. That proposition constructor says
            that, for the same pointer value, a pointer
            that references a larger type provides \textit{less} information 
            than a smaller type. The reason for this is that we want to be able
            to dereference a pointer as close to its bound as possible; a smaller
            type \verb|size| gives us more freedom to dereference anywhere between
            a \verb|base| and \verb|bound - size|.

            \begin{verbatim}
            Inductive leq : alv -> alv -> Prop :=
                | bot_leq_all    :   forall a : alv,   leq bottom a
                | all_leq_top    :   forall a : alv,   leq a top
                | ofst_leq_ofst  :   forall ba p bo s1 s2 : nat, 
                    s1 <= s2 -> leq (fpoffset ba p bo s1) (fpoffset ba p bo s2).
            \end{verbatim}

            Though the $\sqcup$ binary operation is straightforward, I will 
            defer to the \href{https://github.com/parthsarkar17/softboundcets-ese5370/blob/8be2130a53f6cb4e7e283aeb8e46bc1ba478985f/spa.v#L36}{source code}.
            Finally, given this definition of a lattice, I provided an
            implemenation for the following specification. It is long, so I
            again defer to the source file (same as above).

            \begin{verbatim}
            Theorem join_is_monotonic : forall x y z: alv, 
                (leq x y -> leq (join x z) (join y z))
             /\ (leq y z -> leq (join x y) (join x z)).
            \end{verbatim}

            \subsubsection{Montonicity of Flow Rules}

                The next piece of the formalization involved specifying what
                happens to a variable's mapping when it encounters an LLVM
                instruction that assigns that variable. These so-called ``flow rules''
                are also required to be monotonic. Here is my representation of
                LLVM instructions that assign to pointers; aside from  $\phi$-nodes,
                the first three constructors describe the complete list of instructions that compute constant offsets
                from stack-allocated variables:

                \begin{verbatim}
                Inductive llvm_instruction : Type :=
                    | static_alloca (typ_size : nat) (fp_ofst : nat)
                    | bitcast (typ_size : nat) (operand_alv : alv)
                    | gep (typ_size : nat) (ofst : nat) (operand_alv : alv)
                    | other.
                \end{verbatim}


                Then, I described the flow rules in a Rocq function definition,
                written \href{https://github.com/parthsarkar17/softboundcets-ese5370/blob/8be2130a53f6cb4e7e283aeb8e46bc1ba478985f/spa.v#L113}{here}.
                Intuitively, the rule for \verb|static_alloca| provides a variable with a \verb|fpoffset| ALV,
                \verb|bitcast| refines a \verb|fpoffset| value with new type information,
                and the \verb|getelementptr| instruction updates a \verb|fpoffset| with the new \verb|ptr| value
                computed in the pointer arithmetic operation itself. 

                \bigbreak{}

                I provide the theorem and implementation proving the monotonicity
                of flow rules below. I made use of the \verb|auto| Rocq tactic for simpler proofs that are more robust to data structure changes.

                \begin{verbatim}
                Theorem flow_rules_monotonic : 
                    forall insn, leq (flow_rule insn) (flow_rule insn).
                Proof.
                    destruct insn; simpl; auto; 
                    destruct operand_alv; simpl; auto.
                Qed.
                \end{verbatim}
            
    \subsection{LLVM Implementation Details}

        I thought it would be valuable to briefly shed light on some specifics
        regarding my LLVM pass.  I implemented my LLVM dataflow analysis in LLVM 12's standard C++ 
        internal representation. Because the analysis was specific to Softbound+CETS,
        I included the analysis in the same file as the transformation that
        provides Softbound+CETS instrumentation. The analysis is approximately 450
        lines of code (including abstraction behind C++ niceties, like classes).
        The source code is available \href{https://github.com/parthsarkar17/softboundcets-ese5370/blob/009263fa16b6ff7cc22f6afadf9561a9507f60b0/llvm-project/llvm/lib/Transforms/SoftBoundCETS/SoftBoundCETS.cpp#L370}{here}.
    
        \bigbreak{}

        I tried very hard to abide by a direct translation of the Rocq
        description of my flow rules to avoid introducing any compiler bugs. The transfer
        function (that describes how an instruction flows variable to ALV mappings)
        is available \href{https://github.com/parthsarkar17/softboundcets-ese5370/blob/009263fa16b6ff7cc22f6afadf9561a9507f60b0/llvm-project/llvm/lib/Transforms/SoftBoundCETS/SoftBoundCETS.cpp#L601}{here}.
    
        \bigbreak{}

        Next, I actually  use the results of this analysis by passing an object
        representing analysis results to the function that performs the transformation.
    Specifically, I perform the analysis \href{https://github.com/parthsarkar17/softboundcets-ese5370/blob/009263fa16b6ff7cc22f6afadf9561a9507f60b0/llvm-project/llvm/lib/Transforms/SoftBoundCETS/SoftBoundCETS.cpp#L4061}{here},
        and I use the resulting object as an argument to the function \verb|SoftBoundCETSPass::addSpatialChecks|.
        This transformation takes in an instruction and instruments it with a
        call to \verb|spatial_dereference_check()| if it does not early-return.
        So, I use my analysis results precisely 
        \href{https://github.com/parthsarkar17/softboundcets-ese5370/blob/009263fa16b6ff7cc22f6afadf9561a9507f60b0/llvm-project/llvm/lib/Transforms/SoftBoundCETS/SoftBoundCETS.cpp#L3407}{here}, 
        and early-return if the analysis says it's OK to do so!

        \subsection{Correctness, Results \& Evaluation}
            
            The following sections detail the way I evaluated my optimization.
            In the spirit of my Rocq proofs, I felt strongly that the 
            implementation should be thoroughly tested for correctness. I was
            also very curious to see what performance difference my pass would make.

            \subsubsection{Correctness}

                First, I began with hand-written tests. Since my optimization
                attempts to elide dereferences to constant-offset pointers to 
                the stack, I thought of two specific ways my implementation might
                be broken:

                \bigbreak{}

                I first wrote a program containing the following buffer overflow:

                \begin{verbatim}
                    // Test Code A
                    int buf[] = {1, 2, 3};
                    buf[3] = 37;
                \end{verbatim}

                As is expected,  with my optimization enabled, this program fails with a Softbound+CETS-thrown 
                exception. I then made sure my type size propagation was correctly implemented
                by attempting a ``type confusion'' style attack. Here, the pointer
                variable \verb|buf| should only be able to write 3 bytes. However,
                I attempt to cast this pointer to an \verb|(int *)| value, and
                thereby escalate the privilege by 1 byte.

                \begin{verbatim}
                    // Test Code B
                    char buf[] = {1, 2, 3};
                    int *buf_alias = (int *)buf;
                    *buf_alias = 37;
                \end{verbatim}

                Again, my program fails as desired. 
                
                \bigbreak{}

                For a more thorough and representative
                suite of tests, (INSERT [revisited]) provides a way to run
                the Juliet C Test Suite, which is a testing benchmark that
                contains C programs with vulnerabilities. By providing both
                well-formed and ill-formed inputs, compilers protecting against
                memory violations should succeed on valid inputs and abort
                on bad inputs. \href{https://github.com/parthsarkar17/softboundcets-ese5370/tree/009263fa16b6ff7cc22f6afadf9561a9507f60b0/juliet-test-suite-c}{This subdirectory} of my project contains the
                output of the Juliet C Test Suite's stack buffer overflow tests on standard Softbound+CETS (prefixed with \verb|not_spa|)
                and Softbound+CETS with my optimization enabled (prefixed with \verb|spa|).
                I performed a \verb|diff| of the good and bad test case outputs respectively,
                and saw no difference from the original Softbound+CETS implementation.

                \bigbreak{}

                Together with my Rocq implementation proving monotonicity and a
                sufficiently conservative analysis, I believe these two methods of testing
                demonstrate the correctness of my implementation.

            \subsubsection{Results \& Evaluation}

                To evaluate any effects of my optimization, I used the LU-decomposition
                and linear solve benchmark from Section 3. For compiler flags
                \verb|-O0| and \verb|-O3|, I plotted the runtimes of the 
                Softbound+CETS instrumented code both with and without my
                optimization, normalized against the execution time taken by the program compiled with my system's default \verb|clang|
                binary.

                \begin{figure}[H]
                    \centering
                    \includegraphics[width=0.7\textwidth]{/Users/parthsarkar/Desktop/O0-iter4-normd.png}
                    \caption{Normalized Time vs. Compiler Setups (O0, iter=4)}
                \end{figure}

                As demonstrated by the image above, my optimization provides
                an over $2\times$ speedup compared to baseline Softbound+CETS in \verb|-O0|, which was definitely a
                nice surprise. Another surprising result was the fact that baseline
                Softbound+CETS took almost $25\times$ the normal \verb|clang|-compiled binary's running time.
                This is much larger than the average of 140\% cited by (CITATION; revisited).

                \bigbreak{}

                Unfortunately, my optimization seemed to have no effect when each
                binary was compiled with \verb|-O3|:

                \begin{figure}[H]
                    \centering
                    \includegraphics[width=0.7\textwidth]{/Users/parthsarkar/Desktop/O3-iter16-normd.png}
                    \caption{Normalized Time vs. Compiler Setups (O3, iter=16)}
                \end{figure}

                Here are my suspicions. First, \verb|-O3| introduces a myriad of other optimizations that make it more challenging to
                single out the effect of any one optimization. Second, more importantly, \verb|-O3|
                includes a pass that promotes many variables from living in the stack frame to living
                in registers. Therefore, word-sized values that we would initially have to grab from the stack
                frame under \verb|-O0| (each of which previously required a Softbound+CETS dereference check
                that my pass had the opportunity to optimize away)
                can now be retrieved \textit{for free} because they already exist in architectural registers.
                These word-sized values in the stack were prime targets for my pass because
                pointers to these values were all constant offsets from the frame pointer.
                I suspect that my pass has no added effect here in \verb|-O3| because all of these
                word-sized values have now been promoted into registers, and therefore no longer benefit from
                dereference check elision; in fact, they no longer need dereference checks at all!


                \bigbreak{}

                Even though \verb|-O0| typically means ``no optimizations'', I
                believe that my optimization should always be run at \verb|-O0|,
                even if you disregard the $2\times$ performance improvement.
                First, the reason people might want \verb|-O0| is for debugging
                purposes. If that is the case, removing extraneous information
                like runtime dereference checks improves program readability.
                Similarly, people might want to avoid \verb|-O3|, opting instead
                for \verb|-O0|, because higher optimization levels usually
                increase code size. In my pass, code size strictly gets smaller
                because it never adds code; it only removes it.
                

\section{Conclusion \& General Takeaways}

    Broadly, I feel that my project was successful in attempting to motivate
    both encryption of metadata variables and the use of targeted,
    correct compile time optimizations to recover performance sacrificed in the 
    pursuit of security. Given a larger timeline, I would have liked to extend
    my work to achieve full, end-to-end encryption of all metadata variables
    stored to memory. To my understanding, this would have provided the guarantee
    of memory safety for \textit{all} binaries, including those only partially
    compiled with Softbound+CETS. While the overhead would have been very large,
    this further motivates the need for more clever compile time optimizations, including
    interprocedural optimizations (instead of intraprocedural, like the one I
    presented here).
    
    \bigbreak{}
    Even if the overhead is too expensive to warrant deployment as real release 
    code, in my opinion, there is great value in using Softbound+CETS as a debugging tool. Many
    bugs in C (even if never exploited maliciously) arise as a result of illegal
    memory accesses, usually leading to cryptic error messages and segmentation faults. Thus, it would be nice to have an active community
    around Softbound+CETS that routinely upgrades the infrastructure to more
    recent LLVM versions and plugs in nicely to \verb|clang|.


    \subsection*{Personal Fulfillment}

    I deeply enjoyed the process of implementing my dataflow pass, researching
    the relevant theory, and formally proving important properties about it.
    The process helped me feel more confident in my implementation as I was
    writing it, and it generally strengthened my knowledge and interest
    in the subject matter. Additionallly, it was rewarding to learn about and 
    work with the LLVM internal representation. It was not as daunting as I 
    expected, and I now feel more confident approaching tasks that might be 
    better suited as an LLVM-level analysis.
\newpage

\section*{Bibliography}

\end{document}


